\chapter{CORRECTNESS PROOFS AND VALIDATION}\label{chap:correctness}

This chapter establishes the formal logical foundation of the proposed algorithms. We distinguish between \textit{Theoretical Guarantees}—mathematical certainties derived from the algorithm's structure—and \textit{Empirical Observations}—performance metrics that characterize the algorithm's behavior in practice.

\section{Correctness of the Log-Probability Transformation}

\begin{theorem}[Order-Preserving Log Transform]\label{thm:log_transform}
The path $P^*$ with minimum total log-cost is identical to the path with maximum total probability.
\end{theorem}

\begin{proof}
The negative logarithm $f(x)=-\log(x)$ is a strictly monotone decreasing function on the interval $(0,1]$. For any two paths $P_1, P_2$, the transformation preserves the relative ordering of path costs such that the path with the highest cumulative product of probabilities maps to the path with the lowest cumulative sum of negative log-probabilities~\citep{manning1999nlp}. Since $p(u,v) \in (0,1]$, all edge weights $w(u,v) \geq 0$, satisfying the non-negativity constraint required for Dijkstra's algorithm.
\end{proof}

\section{Correctness of Baseline Dijkstra}

\begin{theorem}[Baseline Optimality]\label{thm:baseline_optimal}
Upon termination, $\mathit{dist}[t]$ equals the true shortest-path distance $\delta(s,t)$ in the log-transformed graph.
\end{theorem}

\begin{proof}
We rely on the standard inductive proof for Dijkstra’s algorithm~\citep{cormen2009}. The algorithm maintains the invariant that for every node $u$ extracted from the priority queue, $\mathit{dist}[u] = \delta(s,u)$. Because all edge weights $w(u,v) \geq 0$, the extraction of a node $u$ with the minimum tentative distance guarantees that no shorter path can exist through nodes currently in the queue or yet to be discovered. 
\end{proof}

\section{Correctness of Probability-Pruned Dijkstra}

\subsection{Convergence at \texorpdfstring{$\tau = 0$}{tau = 0}}

\begin{theorem}[Convergence]\label{thm:convergence}
As $\tau \to 0^+$, Algorithm~2 converges to Algorithm~1.
\end{theorem}

\begin{proof}
As $\tau \to 0^+$, we have $T = -\log(\tau) \to +\infty$. Setting $\tau = 0$ is equivalent to $T = \mathtt{math.inf}$ in the implementation. The pruning condition $\mathit{new\_cost} > T$ is never satisfied since all finite costs are $< +\infty$. Therefore no edge is pruned, and every code path of Algorithm~2 behaves identically to Algorithm~1. Correctness follows directly from Theorem~\ref{thm:baseline_optimal}.
\end{proof}

\subsection{Conditional Optimality for \texorpdfstring{$\tau > 0$}{tau > 0}}


\begin{theorem}[Conditional Optimality]\label{thm:conditional}
If the optimal path $P^*$ has a total probability $\Pi^* \geq \tau$, Algorithm 2 returns $P^*$ with an optimality gap $\Delta = 0\%$.
\end{theorem}

\begin{proof}
Let $C^* = \sum_{(u,v)\in P^*} -\log\,p(u,v)$ be the optimal log-cost and $T = -\log(\tau)$ the pruning threshold. By hypothesis $\Pi^* \geq \tau$, so $C^* = -\log(\Pi^*) \leq -\log(\tau) = T$. Consider any prefix $(s = v_0, v_1, \ldots, v_j)$ of $P^*$. Because Dijkstra settles nodes in non-decreasing cost order and $P^*$ is optimal, the cost of each prefix satisfies $\sum_{i=0}^{j-1} w(v_i, v_{i+1}) \leq C^* \leq T$. Therefore no edge along $P^*$ triggers the pruning condition $\mathit{new\_cost} > T$, and $P^*$ remains in the search space. Since the pruned algorithm preserves the priority-queue ordering and settled-node invariant of Theorem~\ref{thm:baseline_optimal} for all unpruned edges, it settles $t$ with cost $C^*$. Hence $\Delta = 0\%$.
\end{proof}

\subsection{The All-or-Nothing Property}

Unlike traditional approximation algorithms that may return a sub-optimal path within an $\epsilon$-bound, the proposed pruning method prioritizes \textbf{Optimality over Reachability}.

\begin{theorem}[All-or-Nothing]\label{thm:all_or_nothing}
Algorithm 2 either returns the globally optimal path or reports no path found. It never returns a sub-optimal path.
\end{theorem}

\begin{proof}
Suppose Algorithm~2 returns a path $P'$ from $s$ to $t$. Since the pruning check ($\mathit{new\_cost} > T$) only skips edges---it does not alter edge weights or the priority queue ordering---the settled-node invariant from Theorem~\ref{thm:baseline_optimal} holds for all nodes explored by the pruned algorithm. Therefore, among all paths that survive the threshold, $P'$ minimizes log-cost. If $P^*$ also survives (Theorem~\ref{thm:conditional}), then $P' = P^*$. If $P^*$ is pruned, then $t$ is unreachable in the pruned subgraph and no path is returned.

To make this explicit, let $C^*$ be the baseline optimal cost and $T=-\log\tau$. If $P^*$ is pruned, then $C^*>T$. For any alternative $s\to t$ path $Q$, baseline optimality implies $C(Q)\ge C^*>T$. Hence no $s\to t$ path can satisfy the pruning constraint, so $t$ is unreachable in the pruned graph. Therefore Algorithm~2 cannot return a different (sub-optimal) path when $P^*$ is pruned.

In either case, the algorithm never returns a sub-optimal path: the gap $\Delta$ is exactly 0\% whenever a path is found.
\end{proof}

\begin{corollary}[Returned Path Is Optimal]
If Algorithm~2 returns a path, that path is the globally optimal shortest path.
\end{corollary}

\section{Empirical Validation and Reachability}

Because we prioritize optimality, the primary ``cost'' of pruning is the potential failure to find a path. We define \textbf{Reachability ($R$)} as the core empirical metric to evaluate this trade-off.

\begin{definition}[Reachability]
The Reachability $R$ is the ratio of successful path-finding by the pruned algorithm compared to the baseline:
\\[-1em]
\[ R = \frac{\text{Count}(\text{Path Found}_{\text{Pruned}})}{\text{Count}(\text{Path Found}_{\text{Baseline}})} \]
\end{definition}



\subsection{Summary of Results}

The properties were verified across 3,000 experimental configurations. As shown in Table~\ref{table:correctness_summary}, the optimality gap remained $0.0000\%$ in all cases where a path was returned, while Reachability varied as a function of $\tau$.

\begin{table}[htbp]
\centering
\renewcommand{\arraystretch}{1.3}
\begin{tabular}{@{}lll@{}}
\toprule
\textbf{Property} & \textbf{Type} & \textbf{Validation Result} \\
\midrule
Order-Preserving & Theoretical & Verified ($<10^{-12}$ precision error) \\
All-or-Nothing & Theoretical & $\Delta = 0.0000\%$ in 100\% of found paths \\
Convergence & Theoretical & Identical output at $\tau = 0$ \\
Reachability ($R$) & \textbf{Empirical} & $R$ decreases monotonically as $\tau$ increases \\
Speedup & \textbf{Empirical} & Up to $18{,}882\times$ on real-world data \\
\bottomrule
\end{tabular}
\caption{Summary of Correctness and Performance Validation.}
\label{table:correctness_summary}
\end{table}

\section{Separation of Claims and Limitations}

\subsection{Proven Theoretical Claims}
The following are formally guaranteed by Theorems~\ref{thm:log_transform} through~\ref{thm:all_or_nothing}:
\begin{itemize}
    \item The algorithm is optimality-preserving (it never returns a "wrong" path).
    \item The log-transformation is mathematically sound for probability maximization.
\end{itemize}

\subsection{Empirical Observations (Expected Outcomes)}
The following depend on graph topology and data distributions and are not formally guaranteed:
\begin{itemize}
    \item \textbf{Search-Space Reduction:} The actual number of edges relaxed depends on the specific $\tau$ and the concentration of high-probability paths.
    \item \textbf{Optimal Threshold $\tau_c$:} Identifying a threshold that maximizes speedup while maintaining $R \approx 1$ is an empirical task performed via a sweep (Chapter~\ref{chap:results}).
\end{itemize}

The following statements are mathematical theorems with formal proofs, independent of experimental data:

\begin{enumerate}
\item[(T1)] The log-probability transformation $w(u,v) = -\log p(u,v)$ preserves path ordering (Theorem~\ref{thm:log_transform}).
\item[(T2)] Baseline Dijkstra returns the globally optimal path in $O(|E|\log|V|)$ time (Theorem~\ref{thm:baseline_optimal}).
\item[(T3)] When $\tau = 0$, Algorithm~2 is identical to Algorithm~1 (Theorem~\ref{thm:convergence}).
\item[(T4)] When the optimal path's probability satisfies $\Pi^* \geq \tau$, Algorithm~2 returns it with zero gap (Theorem~\ref{thm:conditional}).
\item[(T5)] Algorithm~2 never returns a sub-optimal path; it either returns the global optimum or reports no path (Theorem~\ref{thm:all_or_nothing}).
\end{enumerate}

\subsection{Empirically Validated Findings}

The following results are supported by experimental data but \textbf{not formally guaranteed} for arbitrary graphs:

\begin{enumerate}
\item[(E1)] All-rows wall-clock speedups range from $3.3\times$ ($\tau=0.001$) to $738\times$ ($\tau=0.5$) on synthetic graphs, though these figures are dominated by runs where the pruned algorithm terminated without finding a path. Restricting to path-found configurations, wall-clock speedups are $1.6\times$ to $7.2\times$ (Table~\ref{table:speedup_tau}, Chapter~\ref{chap:results}).
\item[(E2)] Speedup scales monotonically with graph size $|V|$ at fixed $\tau$ (Table~\ref{table:speedup_size}).
\item[(E3)] Erd\H{o}s--R\'{e}nyi graphs exhibit higher speedups than layered graphs at identical $\tau$ (Table~\ref{table:speedup_graph_type}).
\item[(E4)] Critical $\tau_c$ exists in a narrow empirical range where speedup peaks without loss of reachability (Chapter~\ref{chap:results}, Section~\ref{sec:speedup_vs_tau}).
\item[(E5)] On real-world datasets under fixed $\tau$, all-rows speedups reach $18{,}882\times$ (RetailRocket item-level) and $2{,}309\times$ (RecSys~2015), reflecting rapid termination without path recovery. Under adaptive $\tau$, found-only speedups are up to $5.8\times$ (Chapter~\ref{chap:results}, Section~\ref{sec:real_results}).
\item[(E6)] Optimality gap remains $0.0000\%$ across 706 path-found cases in combined synthetic and real experiments (all tested configurations).
\end{enumerate}

\subsection{Acknowledged Limitations}

The following limitations are explicitly recognized:

\begin{enumerate}
\item[(L1)] No formal approximation guarantee exists for $\tau > 0$. While we prove convergence (\textit{T3, T4}), the pruned algorithm's behavior for a specific $\tau$ on an arbitrary graph is not formally characterized. Speedup and reachability trade-offs are purely empirical.
\item[(L2)] Critical $\tau_c$ is defined empirically; we provide no theoretical method to predict it a priori for a given graph without experimental sweep.
\item[(L3)] Results assume edge weights are non-negative (the log-probability transformation guarantees this, by construction). The framework does not extend to negative weights or negative-cycle detection.
\item[(L4)] Real-world validation is limited to two e-commerce clickstream datasets. Generalization to other domains (social networks, transportation, etc.) is not formally justified and would require further empirical testing.
\end{enumerate}
